\subsection{(Optional) Textual Generate}
To enhance the expressiveness and personality of generated stickers, we implemented an optional text overlay pipeline that adds short, creative captions directly onto the final sticker images. This subsection describes the three key components: the limitation of diffusion-based text generation, VLM-powered caption generation, and style-aware text rendering.

\subsubsection{\textbf{Limitation of Direct Text Generation}}
Stable Diffusion 1.5, the base model used in our pipeline, cannot reliably generate legible text within images. When text is included in the generation prompt (e.g., ``a sticker of a cat with the text `Cute Cat'\,''), the model produces garbled or distorted characters. This is a well-known limitation of diffusion models that lack dedicated text-rendering modules. Therefore, we adopt a post-processing approach: captions are overlaid onto the generated sticker images using PIL (Python Imaging Library) after the image generation stage.

\subsubsection{\textbf{VLM-Based Caption Generation}}
Rather than manually writing captions for each sticker, we employ a Vision-Language Model (VLM) to automatically generate contextually appropriate captions by analyzing the sticker image content. Specifically, we use GPT-4o-mini via the POE API with a carefully engineered system prompt.

The VLM is instructed to output exactly two words in the format of \textbf{adjective + noun} (e.g., ``Sleepy Whiskers'', ``Brave Bear'', ``Chill Penguin''). The prompt includes diversity constraints to prevent repetitive nouns---for instance, for cat stickers, the model is encouraged to rotate among alternatives such as Cat, Kitty, Whiskers, Furball, Meowster, and Fluffball. This produces captions that are both descriptive and varied across the sticker set.

The generated captions are cached in a JSON file (\texttt{text\_captions\_vlm.json}) for reproducibility and reuse, avoiding redundant API calls during subsequent pipeline runs.

\subsubsection{\textbf{Style-Aware Text Rendering}}
A key challenge in text overlay is ensuring readability across diverse sticker backgrounds. A fixed text color (e.g., white) may blend into light-colored regions, making the caption illegible. To address this, we implement a style-aware rendering system with three components:

\textbf{Dominant Color Extraction.} For each sticker image, we extract the most prominent saturated color by downsampling the image to 50$\times$50 pixels, filtering out extreme brightness and low-saturation pixels, and selecting the most frequent quantized color. This dominant color is used to inform the text color scheme.

\textbf{Local Background Brightness Sampling.} Before placing text, we sample the average brightness of the image region directly behind the text bounding box. This local measurement is more accurate than global image statistics, as sticker backgrounds often vary significantly between the center and edges.

\textbf{Contrast Adjustment.} We compare the intended text fill color against the local background brightness. If the contrast is insufficient (difference $< 60$ on a 0--255 scale), the fill and stroke colors are automatically swapped to a high-contrast alternative, ensuring the caption remains legible regardless of background.

Each of the three sticker styles uses a distinct rendering approach:
\begin{itemize}
    \item \textbf{Flat style}: Arial Bold font with the dominant image color as fill, white stroke, and a subtle drop shadow for depth.
    \item \textbf{Cartoon style}: Comic Sans Bold font with white fill and a thick colored outline (4px stroke), mimicking comic book lettering.
    \item \textbf{Watercolor style}: Segoe Script (handwritten) font with a soft Gaussian glow behind the text, complementing the painterly aesthetic.
\end{itemize}

Fig.~\ref{fig:text_3styles} shows an example of the three rendering styles applied to the same animal sticker, demonstrating how each style produces a visually distinct yet readable caption.

\begin{figure}[htbp]
    \centering
    \includegraphics[width=0.95\columnwidth]{Figures/text_overlay_3styles_animal_1.png}
    \caption{Text overlay comparison across three sticker styles. Left: Flat (dominant color + shadow). Center: Cartoon (white fill + colored outline). Right: Watercolor (handwritten font + glow).}
    \label{fig:text_3styles}
\end{figure}

Fig.~\ref{fig:text_grid} presents a broader view with multiple animals across all three styles, illustrating the consistency and adaptability of the text rendering pipeline.

\begin{figure}[htbp]
    \centering
    \includegraphics[width=0.95\columnwidth]{Figures/text_overlay_grid.png}
    \caption{Text overlay results for multiple stickers across three styles. Each row represents a different animal; each column represents a style. Captions are generated by VLM and rendered with style-specific fonts and contrast-aware colors.}
    \label{fig:text_grid}
\end{figure}

This style-matched approach ensures that the text overlay feels visually integrated with each sticker style rather than appearing as an afterthought. The entire pipeline processes all stickers in batch, producing final output images with consistent quality.
